

\begin{derivation}[从\eqnrefs{eq:bubble-imag-phase-space-dp8}到\eqnrefs{eq:two-body-threshold-beta-dp8}]

正文\eqnrefs{eq:bubble-imag-phase-space-dp8} 表明虚部只有在两个内部粒子同时满足正能质量壳和能量守恒时才存在。对等质量情形进入质心系，能量 $\delta$ 函数固定唯一的径向动量 $k_*$；解出这个根并把它写成相对论速度因子，就得到正文\eqnrefs{eq:two-body-threshold-beta-dp8}。

在质心系 $\mathbf P=0$ 且 $m_1=m_2=m$ 时，
\begin{equation*}
 E_1=E_2=E_k=\sqrt{m^2c^4+c^2k^2},
 \qquad
 \sqrt{s}\,c=2E_k.
\end{equation*}
平方后有
\begin{align*}
 sc^2
 &=4m^2c^4+4c^2k_*^2,\\
 k_*^2
 &=\frac{s-4m^2c^2}{4}.
\end{align*}
因此只有 $s\ge4m^2c^2$ 时 $k_*$ 才为实数。再用 $E_{\rm cm}=\sqrt{s}\,c$，
\begin{align*}
 \beta_*(s)
 &=\frac{ck_*}{E_{\rm cm}/2}
 =\frac{2ck_*}{\sqrt{s}\,c}\\
 &=\sqrt{1-\frac{4m^2c^2}{s}},
\end{align*}
即正文\eqnrefs{eq:two-body-threshold-beta-dp8}。为把阈值因子的整体归一化也固定下来，定义正文\eqnrefs{eq:bubble-imag-phase-space-dp8} 中的径向相空间核
\begin{equation*}
 \rho_2(s)
 \equiv
 \int\frac{\dd^3\mathbf k}{(2\pi\hbar)^3}
 \frac{c^2}{4E_k^2}
 \delta\!\left(\sqrt{s}\,c-2E_k\right).
\end{equation*}
利用
\begin{equation*}
 \left.\frac{\dd(2E_k)}{\dd k}\right|_{k_*}
 =\frac{2c^2k_*}{E_*},
 \qquad E_*\equiv E_{k_*}=\frac{\sqrt{s}\,c}{2},
\end{equation*}
径向 $\delta$ 积分严格给出
\begin{align*}
 \rho_2(s)
 &=\frac{4\pi k_*^2}{(2\pi\hbar)^3}
 \frac{c^2}{4E_*^2}
 \frac{E_*}{2c^2k_*}\\
 &=\frac{k_*}{16\pi^2\hbar^3E_*}
 =\frac{\beta_*(s)}{16\pi^2\hbar^3c}.
\end{align*}
因此连续谱在阈值处按确定的线性因子 $\beta_*(s)$ 从零打开。
\end{derivation}

\begin{derivation}[从\eqnrefs{eq:optical-unitarity-r68}到\eqnrefs{eq:perturbative-optical-theorem-dp68}]

正文\eqnrefs{eq:optical-unitarity-r68} 是精确散射算符的幺正性关系。把跃迁算符按耦合阶数展开，并只比较最低出现虚部的二阶项，就能把圈振幅的虚部写成一阶真实跃迁振幅的模平方之和，从而得到正文\eqnrefs{eq:perturbative-optical-theorem-dp68}。

写
\begin{equation*}
 \hat T_S=\hat T^{(1)}+\hat T^{(2)}+\hat T^{(3)}+\cdots .
\end{equation*}
代入
$-\ii(\hat T_S-\hat T_S^\dagger)=\hat T_S^\dagger\hat T_S$，二阶部分为
\begin{equation*}
 -\ii\bigl(\hat T^{(2)}-\hat T^{(2)\dagger}\bigr)
 =\hat T^{(1)\dagger}\hat T^{(1)}.
\end{equation*}
在同一初态 $|i\rangle$ 两侧取矩阵元，并在右端插入在壳中间态完备关系 $\sum_X|X\rangle\langle X|=\mathbb I$，得到
\begin{align*}
 -\ii\left[
 \langle i|\hat T^{(2)}|i\rangle
 -\langle i|\hat T^{(2)}|i\rangle^*
 \right]
 &=\sum_X
 \langle i|\hat T^{(1)\dagger}|X\rangle
 \langle X|\hat T^{(1)}|i\rangle\\
 &=2\operatorname{Im}\langle i|\hat T^{(2)}|i\rangle,
\end{align*}
这就是正文\eqnrefs{eq:perturbative-optical-theorem-dp68}。图形语言中的割线规则只是把右端的在壳中间态完备和翻译成内部传播子的 $\delta_+$ 替换。
\end{derivation}
